\subsection{Three-object stars at the nineteen-arrow frontier}

\begin{proposition}[Proved: a small central corner through nineteen arrows]
\label{CP2:A-star-small-center-nineteen}
A connected three-object source or sink star with at most nineteen
arrows and central corner of order at most four cannot have homotopy
type $\CP^2$.
\end{proposition}
\begin{proof}
Use opposition to make the center a source, and write the corner
orders as $a,b,c$ and the two nonempty stalk sizes as $s,t$.
The small-monoid comparison gives $BM\simeq BG$ for the minimal
finite group of the central monoid. If $G$ is trivial, the actual
union of the two edge nerves over $BM$ is homotopy equivalent to
the wedge of the two edge spaces. A wedge with projective-plane
type has one edge of projective-plane type: its fundamental group
is a free factor of the trivial group, integral homology is a direct
summand, and the unique primitive second square must occur in that
summand. Each edge has at most $19-2=17$ arrows. This contradicts
Theorem~\ref{CP2:A-global-lower-nineteen}.

If $G$ is nontrivial, the minimal ideal of $M$ is itself a group.
Indeed a Rees ideal with a nontrivial group and at least two rows
or columns already has at least four elements, with the distinct
global identity outside it, exceeding $a\le4$. Averaging over the
group ideal makes the trivial left and right rational modules
projective. The actual exact-orbit functor replacing the center by
the trivial monoid is a rational homology and cohomology-ring
equivalence. Its target is a wedge of two ordinary act cones, with
leaf monoids of orders $b,c$ and state counts at most $s,t$.

A rational projective-plane ring on this wedge would force one
summand to have exactly that rational homology and ring: the
second and fourth cohomology are each one-dimensional and the
nonzero square belongs to one summand, while all other positive
homology groups of both summands vanish. If the relevant leaf
order is at most eleven, the bounded-eleven act theorem excludes
it. Since now $a\ge2$ and the other leaf and stalk each have at
least one arrow, a leaf of order twelve has at most three states.
Proposition~\ref{CP2:A-twelve-three-state-rational} excludes it.
A leaf of order at least thirteen has at most two states, so the
unrestricted binary-act product obstruction applies. These exhaust
the budget.
\end{proof}

\begin{proposition}[Proved: a five-element coned center with a leaf through nine]
\label{CP2:A-five-center-nine-four}
Let $M$ have order five, let $N$ have order at most nine, and let
$S$ be a nonempty $(N,M)$-biset with at most four states. The
cofiber of $B(N,S,M)\longrightarrow BN$ cannot have homotopy
type $\CP^2$.
\end{proposition}
\begin{proof}
We spell out the budget extension of
Theorem~\ref{CP2:A-star-five-seven-four}. Its preceding lemma
\ref{CP2:A-coned-four-center-P5} handles $M=P_5$ for arbitrary
$N$. For every other five-element monoid, $BM$ is equivalent to
its minimal group. If that group is trivial, coning the central
corner does not change the two-object collage up to homotopy.
That collage has at most $5+9+4=18$ arrows and is excluded by
Theorem~\ref{CP2:A-global-lower-nineteen}.

If the minimal ideal has one row, exact rational orbit comparison
identifies the cofiber with the ordinary $N$-act cone on $S/M$ in
rational homology and ring, through the actual orbit functor.
The proposed target supplies bounded homology; the bounded-eleven
act theorem excludes its square. The only remaining minimal-group
possibility under the order-five budget is
$M=(\{0,1\}\times C_2)^1$, with
$(i,g)(j,h)=(i,gh)$.

Put $U=SK$. If $U$ consists of one $C_2$ orbit, the commuting
$N$ action permutes it. The finite-cover transfer for
$B(N,U)\to BN$, factoring through the bridge, makes every
second class of the cofiber relative to $BN$, and its square is
zero. A unique singleton core orbit is a common fixed state and
gives the same conclusion by a section. If $S=U$, every right
orbit module is projective over $\mathbb QM$: it is a summand
of copies of $e\mathbb QM\cong\mathbb QC_2$. The actual orbit
augmentation therefore gives the same ordinary-cone rational
comparison as above. The bound $|S|\le4$ leaves only the case
in which every core orbit is a singleton.

In that case the right action factors through
$\overline M=\{0,1\}^{\mathrm{left},1}$, and choosing the
identity group coordinate gives a literal monoid section. The
quotient star is an actual retract with two fewer arrows. The
kernel of $\mathbb QM\to\mathbb Q\overline M$ is
$\mathbb QM p$, where
$p=((0,1)-(0,-1))/2$ is idempotent, and $\mathbb Q[S]p=0$.
The resulting length-one projective comparison shows that this
actual retract is a rational homology equivalence. Its integral
homology is a split summand of that of $\CP^2$; the rational
second isomorphism and natural primitive square force its second
and fourth groups to be $\mathbb Z$ with a primitive square.
Its fundamental group is a retract of the trivial group. Integral
recognition consequently gives actual projective-plane type.
The original coned star has at most $5+9+4+2=20$ arrows, so this
retract has at most eighteen, again a contradiction.
\end{proof}

\begin{corollary}[Proved: two additional cardinality cuts for nineteen-arrow stars]
\label{CP2:A-nineteen-star-pendant-and-coned}
In a nineteen-arrow source star, either a nontrivial leaf on a
singleton stalk, or a trivial leaf on a singleton stalk whose
other stalk has at most three arrows, excludes projective-plane
type. The same holds after opposition.
\end{corollary}
\begin{proof}
For the first alternative, the arbitrary singleton-pendant
retract theorem replaces the leaf corner by its identity and
preserves actual projective-plane type. At least one arrow is
lost, contradicting the established eighteen-arrow exclusion.
For the second alternative the space is the coned bridge to the
other leaf. That leaf has order at most fifteen, since the
center, the two stalks and the trivial leaf each contribute at
least one arrow. Corollary~\ref{CP2:A-coned-three-leaf-bounds}
therefore excludes actual projective-plane type.
\end{proof}
